\documentclass[12pt]{article}
\usepackage[T1]{fontenc}
\usepackage[utf8]{inputenc}


\usepackage{geometry}
\geometry{a4paper, margin=2.5cm}
\usepackage{setspace}
\usepackage{natbib}
\usepackage{hyperref}
\usepackage{url}
\usepackage{booktabs}
\usepackage{csquotes}


\setlength{\parindent}{1em}
\setlength{\parskip}{0.5em}

\title{Reconfiguration of Inquiry in LLM Environments:\\
Origin Opacity and the Conditions of Human Responsibility}

\author{}

\date{}

\begin{document}

\maketitle

\begin{abstract}
As large language models (LLMs) become infrastructural, the transformations they bring about can no longer be adequately grasped through the classical framework of technology altering the conditions of human activity. Wiener, Ellul, and Arendt already demonstrated that mechanization, automation, and cybernetics reconfigure the conditions of human responsibility, the meaning of action, and the structure of judgment. This paper inherits that historical trajectory while arguing that LLM environments extend such transformation into a qualitatively distinct layer: the pre-reflective generative structure of inquiry itself.

The central question of this paper is not how closely AI comes to resemble human beings. The problem, rather, lies in the possibility that LLM environments alter the origin conditions and generative modalities of human inquiry. Drawing on Gibson's ecological theory of affordances, Polanyi's structure of tacit knowledge, and Merleau-Ponty's phenomenology, this paper argues that probabilistic responses and semantic smoothing in LLM environments can compress the historicity, ambiguity, and indeterminacy that inquiry inherently carries, and that their repetition can reconfigure the field of the questionable itself. As this process accumulates, inquiry continues to be experienced as the subject's own---while the generative conditions of inquiry become increasingly opaque to the subject.

Contemporary AI research is advancing toward augmenting LLMs with neural signals, embodiment, and affective responsiveness, intensifying anthropomorphic tendencies. Yet such developments do not confer upon AI the capacity to bear terminal responsibility. On the contrary, the more the origin of inquiry becomes opaque, the deeper the asymmetry in which final responsibility remains with the human being.

This paper therefore reformulates irreversibility, fixation of responsibility, and terminal human responsibility not as supplementary ethical additions after the fact, but as conditions for technological progress to remain humanly bearable. Its argument is not opposition to technological progress itself, but rather the articulation of the minimum conditions required to preserve the conditions of inquiry and responsibility in LLM environments---updating for the LLM era the classical philosophical claim that technology transforms the human condition.
\end{abstract}

\noindent\textbf{Keywords:} Large Language Models $\cdot$ Irreversibility $\cdot$ Terminal Responsibility $\cdot$ Structural Transformation of Inquiry $\cdot$ Origin Opacity $\cdot$ Philosophy of Technology $\cdot$ Semantic Smoothing

\section{Introduction}

In current discourse on AI, vocabularies of responsibility, ethics, oversight, and explainability have come into wide use. Yet much of this discourse presupposes the advance and social deployment of AI, and asks what adjustments or correctives should be applied within that frame. The problem is frequently formulated as ``how to proceed safely,'' while questions such as ``where ought we to stop'' and ``what must humans still retain'' tend to recede into the background. The purpose of this paper is to describe structures already becoming visible in LLM environments and to examine, as a matter of possibility, what transformations these structures may bring to the conditions of human inquiry and responsibility. In what follows, the term \textit{AI} is used to refer to artificial intelligence systems in a broad sense, while \textit{LLM} refers to specific interaction with conversational language models.

This paper addresses the situation and conditions that arise under the advancing infrastructuralization and anthropomorphization of LLMs. In LLM environments in particular, the problem does not remain confined to the quality of answers or the precision of decision support. LLMs do not merely assist activity and judgment; they can penetrate the pre-reflective conditions themselves from which inquiry arises. What changes in this case is not only what humans do, but potentially the generative conditions of inquiry itself.

The purpose of this paper is to situate this transformation within the genealogy of the history of technology. That technology alters the conditions of human existence is not new. Yet in LLM environments, that transformation may reach the structure of inquiry itself---a fact that indicates the need to reconstruct responsibility theory from a deeper layer. This paper argues such a proposal not as resistance to progress, but as the presentation of conditions under which progress can remain compatible with human responsibility.

The concern of this paper is not to discuss the transformation of the information environment in general terms. In contrast to adjacent discussions that frame the problem as a reconfiguration of the information environment, this paper focuses on a more limited layer: how the pre-reflective generative structure of inquiry itself may undergo transformation.\footnote{On the proximity to Floridi's concept of the ``infosphere,'' see \citet{floridi2014}.}

\section{Technology as Transformer of the Human Condition: A Historical Trajectory}

Since modernity, technology has repeatedly altered the conditions of human activity, judgment, and responsibility. The Industrial Revolution incorporated human skill and discretion into mechanical processes; mechanization increased the efficiency of labor while simultaneously redefining the position of human beings within the productive process. Automation and cybernetics carried the risk of repositioning human beings not as responsible agents of judgment but as elements within complex systems. Wiener warned of the danger of human beings being treated as cogs, while Ellul understood \textit{technique} not as a set of discrete tools but as an environment within which human beings are compelled to define themselves.

Arendt, too, understood automation not merely as an increase in efficiency but as an event that alters the configuration of meaning within human activity. For her, the problem was not that labor becomes lighter, but how such changes transform the structure of human activity as a whole. In this respect, Arendt already saw that technology does not remain external to the human being but enters into the human condition itself.

In this sense, the problem of modern technology lies not only in what it makes possible, but in the fact that it alters the very conditions under which human beings act, judge, and take responsibility.

Subsequent research has continued this trajectory. Studies drawing on extended mind theory and the concept of cognitive niche construction---such as the work of \citet{heersmink2021}---share with this paper a concern for how environments form cognition. Yet where extended cognition addresses the expansion of cognition through environmental coupling, what this paper describes runs in the opposite direction: specifically, the pre-reflective generative structure of inquiry undergoes compression rather than extension. It is the accumulation of this reduction that gives rise to origin opacity. This paper focuses on this more limited layer.\footnote{Coeckelbergh's work on AI ethics and the political and social situatedness of technology \citep{coeckelbergh2020} also makes important contributions, though the layer specific to this paper---intervention in the generative structure of inquiry---is not its primary focus.}

What these discussions primarily addressed was the transformation of the conditions of labor, action, judgment, organization, and responsibility---and how technology repositions human beings within those conditions. Yet in LLM environments, such transformation may reach a still deeper layer: the pre-reflective generative structure of inquiry itself. This is what this paper inquires into.

\section{The Structure of Reception through Embodiment and Anthropomorphization: Toward the Threshold of Inquiry's Transformation}

This chapter does not develop the central argument of this paper. The claim that LLM environments can intervene in the generative structure of inquiry acquires its full reach only under the condition that LLMs are experienced not merely as systems that return answers, but as entities that human beings naturally accept and to whom they are inclined to delegate judgment. This chapter presents the sociotechnical conditions for such reception---the advance of embodiment and anthropomorphization in contemporary AI research---as auxiliary background. The central theoretical development proceeds in Section~4.

What is at issue here is that the transformation arising in LLM environments does not simply repeat this historical trajectory. What Arendt observed in automation was the transformation of the conditions and configuration of meaning within human activity. What this paper addresses, by contrast, is a situation in which the pre-reflective generative structure of inquiry itself can be reconfigured by the responsive environment---a transformation that exceeds mere activity support.

Contemporary AI research is advancing in the direction of augmenting the human-likeness of LLMs and generative AI. NeuroAI research has identified embodiment, language and communication, robotics, and learning as primary points of contact between neuroscience and AI, and the direction of embodied AI occupies an important position in this research landscape \citep{fellous2026,nihbrain2024}. Further, research connecting EEG with LLMs presents frameworks for linking neural signals to affective analysis and inference \citep{neuroai2024}. In short, current research trends are moving toward approaching AI not as a device for mere symbol manipulation, but as an entity equipped with embodiment and neural responsiveness.

This tendency can intensify the anthropomorphization of AI. Human-like responses, affective expression, and the simulation of embodiment can increase user trust and acceptance, and reduce resistance to the delegation of judgment. That is, the more AI is felt to ``understand,'' ``feel,'' and be ``trustworthy,'' the more readily human beings accept its mediation as something natural.

Yet the central asymmetry remains unresolved. The simulation of embodiment, the incorporation of neural signals, and the refinement of affective response can make AI behavior appear human-like, but cannot confer upon it the subjectivity to bear terminal responsibility. Human embodiment is formed through the accumulation of irreversible history---injury, aging, death, the transformation of relationships. AI has no life or death in this sense. Therefore, however refined the simulation of embodiment becomes, it cannot reproduce the irreversible accumulation of history that characterizes human embodiment, nor can it provide AI with the conditions for bearing terminal responsibility.\footnote{EU AI Act Article~14, Article~26(2); Directive~(EU)~2024/2853; \citet{santonideio2018}.}

As a distinct pathway from anthropomorphic reception, it is necessary to address the amplification effect inherent in the response structure of LLM environments. A perceptual structure in which repeated exposure causes what initially felt alien to become familiar---this is already observable in the domain of visual representation. An analogous structure can operate in LLM responses as well. When a particular worldview or discourse is embedded as a system prompt or context, the LLM returns outputs aligned with that vocabulary and framework. When such outputs reach recipients through repeated dialogue, the recipients cannot see the input conditions under which the responses were generated. Responses are readily received as expressions from a neutral intelligence, and the embedded worldview sediments, turn by turn, as self-evident premise. The problem does not lie in malicious intent but in the structural properties of the response architecture of LLM environments. This structure may contribute, at the level of design, to the conditions under which the reconfiguration of the origin conditions of inquiry becomes possible. The theoretical articulation of that process is developed in Section~4.

\section{The Reconfiguration of Inquiry in LLM Environments}

\subsection{Theoretical Grounds for Pre-Reflective Origin Conditions}

\subsubsection{The Structure of Tacit Knowledge: Why Compression Remains Imperceptible}

According to \citet{polanyi1966}, knowing is constituted within a two-layered structure of awareness. Subsidiary awareness supports focal awareness; when inquiry arises, the context, embodiment, and sedimented experience that support its formation operate within the subsidiary dimension. What is crucial is that when one attempts to focalize the subsidiary dimension, its function collapses. When Polanyi states that ``we can know more than we can tell,'' what he points to is precisely this structural asymmetry.

When LLM environments act upon the origin conditions of inquiry, that action proceeds within the subsidiary dimension. The subject continues to feel that they are raising inquiry. Yet when they attempt to focalize where that inquiry comes from, the subsidiary conditions that supported its origin have already undergone transformation. This is the epistemological ground of origin opacity. That compression remains difficult for the subject to perceive is not due to a lack of attention. Action upon the subsidiary dimension is, structurally, incapable of becoming a direct object of observation.

\subsubsection{The Silent Logos: What Is Being Compressed}

In Merleau-Ponty, the ``silent logos'' refers to the proposition that a meaningful openness to the world is established prior to linguistification \citep{merleauponty1968}. Perception is not the passive reception of neutral sensory data but the movement of prior meaning-formation through which the body takes up its relation with the world. The pre-linguistic layer from which inquiry can arise---what this paper calls the ``field of the questionable''---is grounded in the dimension of this silent logos.

Gibson's theory of affordances is useful for showing that inquiry arises within a relation with the environment rather than from the interior of an isolated subject \citep{gibson1979}. Yet what affordances describe are the environmental givens of perception and action. The pre-linguistic and pre-reflective conditions from which inquiry can arise---what can be sensed as inquiry at all---are not contained therein. What Gibson can establish is only that the arising of inquiry is not independent of the environment. To describe the conditions of formation of the field of the questionable---its pre-linguistic ground---the phenomenology of Merleau-Ponty becomes necessary.

When LLMs act upon the field of the questionable, what is being compressed is not the configuration of affordances. What can be reduced is the dimension of the silent logos---the pre-linguistic horizon, formed through irreversible bodily history, within which inquiry can germinate as meaning.

\subsubsection{The Contrast with Chiasm: Asymmetric Reduction in LLM Environments}

Merleau-Ponty's chiasm designates the structure in which the boundary between the perceiving subject and the world mutually interpenetrates \citep{merleauponty1968}. ``The hand that touches is also the hand that is touched''---in this reversibility, subject and world constitute each other while remaining open to each other. The chiasm occurs between a subject bearing irreversible bodily history and the world. What is crucial is that for the chiasm to obtain, both parties must share a common ground: a historicity that is irreversibly, bodily, and existentially assumed. The reference to chiasm in this section is not made in order to apply this concept to the interaction between subject and LLM environment. Rather, it is employed as a contrast to show that the conditions for chiasm are absent in LLM environments. It should be noted that the reference to chiasm in this paper is not a strict reproduction of Merleau-Ponty's original concept, but a limited reconstruction employed to show contrastively what is absent in LLM interaction.

Dialogue with an LLM appears, in formal terms, as bidirectional interaction. The subject receives a probabilistic response, forms the next inquiry, and that inquiry in turn conditions the LLM's next response. Yet here a decisive difference obtains. LLMs possess a record in the sense of logged history, updates, and state transitions. What LLMs do not possess is a historicity that is irreversibly, bodily, and existentially assumed---a history that sediments through injury, aging, the transformation of relationships, and loss. What is at issue in the chiasm is not the absence of the former but of the latter. LLMs do not possess common ground in this sense.

When bidirectional interaction proceeds under this absence, what occurs? The structure of probabilistic responses can induce the subject to complete, unconsciously, what is not visible. This completion remains only on the subject's side. Yet the LLM does not accumulate that completion as bodily history. The completion is internalized by the subject as input for the next cycle, but it does not necessarily remain preserved as the origin of inquiry. Within the nested structure of repetition, the completion itself can become the condition for further compression.

The boundary does not disappear. Yet as the subject enters more deeply into repeated interaction with the LLM environment, that boundary becomes increasingly imperceptible to the subject. The problem is that this imperceptibility itself is not visible to the subject. The Polanyian structure---that what operates in the subsidiary dimension cannot be focalized---functions here once again. The opacity of the boundary is not due to lack of attention; it can arise structurally.

In this structure, formal bidirectionality and effective asymmetry are compatible. Within the internal relation between the subject and the field of inquiry, mutual activation occurs. Yet the effect of interaction with LLMs is asymmetric. It is only on the side of the subject---who possesses a pre-linguistic horizon---that the compression of multiplicity, historicity, and indeterminacy accumulates. This is the structural ground of compression-reconfiguration.

The theoretical grounds established in the preceding section put in place the framework for describing the reconfiguration of inquiry. The field of the questionable is grounded in the pre-linguistic and pre-reflective layer, and the subsidiary conditions that support its formation are not available as objects of direct observation. With this structure as premise, the following describes the dynamic process of the reconfiguration of inquiry in LLM environments.

What proves useful here is Gibson's ecological theory of affordances. It must be noted, however, that what Gibson can establish is only that the arising of inquiry is not independent of its environmental relation. The pre-linguistic conditions of formation of the field of the questionable have been discussed in the preceding section; this section takes that layer as premise and describes the dynamic process that occurs upon it.

What was observed in the preceding chapter---the advance of embodiment and anthropomorphization---indicates the current situation of AI research. Yet the central problem of this paper does not lie in this augmentation of human-likeness itself, but in the influence LLM environments may exert upon the generative structure of inquiry. Technology has always assisted judgment. Yet LLMs do not merely return answers through probabilistic responses; through their repetition, they can condition the directionality of generation itself---that is, what kinds of inquiry are liable to arise next.

From Gibson's ecological perspective, inquiry arises not solely from within the interior of an isolated subject but within the field of possibility opened by the environment \citep{gibson1979}. This perspective is useful here because it makes it possible to understand inquiry not as pure spontaneity residing within the subject's interior, but as something generated within a relation with the environment. In LLM environments, this field of possibility is not static but changes in the course of interaction. The subject receives a probabilistic response from the LLM and returns a further inquiry. In the course of this repetition, a sense of what the subject is liable to receive as self-evident inquiry, and what tends to recede into the background as already sufficiently processed, is gradually formed. What operates here is \textit{semantic smoothing}. Semantic smoothing designates the operation by which probabilistic responses, at the level of content, order ambiguous and indeterminate inquiry into a form that appears more manageable, smoother, and already processed.

What is at issue in LLM environments is therefore not only the accuracy of answers. The problem lies in the fact that the field of the questionable itself undergoes transformation, and that this transformation is not sufficiently visible to the subject. What arises here is the reconfiguration of the pre-reflective origin conditions of inquiry---a transformation that reaches a layer deeper than changes in external conditions of activity. The conditions of formation of this layer have been discussed in the preceding section; what is described here is how the structure operating in the subsidiary dimension is reconfigured through repeated interaction.\footnote{The term ``questionable'' refers to a concept grounded in the dimension of the silent logos discussed in Section~4.1.2: not something already formulated as inquiry, but the pre-linguistic layer, formed through irreversible bodily history, from which inquiry can arise. When LLMs act upon this layer, they are compressing what they are structurally incapable of possessing themselves.}

This transformation does not occur as a single event but can proceed cumulatively through the repetition of inquiry, probabilistic response, and re-inquiry. This repetition is not simple reiteration; it has a nested structure in which the inquiry and probabilistic response of the preceding stage are internalized as the conditions of the next stage's inquiry. What arises in the course of this repetition is \textit{compression}. Compression designates the operation by which the historicity, ambiguity, and indeterminacy that inquiry inherently carries are reduced through probabilistic responses and semantic smoothing.

When compression accumulates in the course of repetition, the field of the questionable itself is reconfigured. This paper calls this \textit{compression-reconfiguration}. Compression-reconfiguration designates the situation in which, when the compression of the historicity, ambiguity, and indeterminacy of inquiry accumulates, the very structure of what is liable to arise as inquiry and what tends to recede as already processed undergoes transformation.

In this sense, the reconfiguration of inquiry in LLM environments is not a matter of cognitive assistance or retrieval efficiency. It concerns the reconfiguration of the very preconditions that underlie the generation of inquiry---what is asked, how it is asked, and what is processed as already having been asked. The field of inquiry and the conditions of the subject mutually activate and change. The reconfiguration of inquiry and the reconfiguration of the subject are not two separate events but two aspects of a dynamic process that proceeds simultaneously.

\section{The Loss of Origin Transparency}

This paper formulates the following process not as an empirically demonstrated universal fact, but as a possibility of structural transformation observable in currently developing LLM environments.

The subject still feels that they are raising inquiry. Yet how that inquiry comes about---from what history, embodiment, relationships, and sedimented experience it arises---may become increasingly difficult to grasp under highly developed LLM environments. LLM environments form the direction of inquiry through the repetition of probabilistic responses, and can cause that formation to be received as if it were unmediated, self-evident inquiry. The subject does not consciously judge each such change individually. Rather, within the flow of interaction, by unconsciously completing and accepting what is not visible, they may lose sight of the transformation of the origin conditions of inquiry---as if it had always been that way.

Origin transparency is lost in this sense. Inquiry continues to be received as one's own inquiry, but that sense of ``mine'' no longer guarantees origin transparency. On the contrary, it is precisely because the sense that inquiry is one's own can itself be maintained through probabilistic responses, semantic smoothing, and unconscious completion that the transformation of origin conditions becomes increasingly imperceptible.

What is important here is that this opacity does not immediately signify exemption from responsibility. The more the origin of inquiry becomes opaque---that is, the closer the subject approaches a state in which they perceive their responses to inquiry to be aligned with those of the LLM---the stronger the temptation to delegate the responsible subject to AI may become. The more the origin of inquiry becomes opaque through repeated smooth dialogue, the more the subject tends to feel that understanding has been achieved and trust has been formed in dialogue with AI. This sense can operate in the direction of causing the subject to accept the delegation of judgment to AI as something natural, while concealing the absence of common ground. Yet that temptation is structurally mistaken. Even as dialogue deepens, the final undertaking remains with the human being. The decrease in origin transparency is not the same as the transfer of the responsible subject. The problem, rather, lies in the asymmetry in which responsibility remains with the human being while the origin conditions of inquiry that support that responsibility become difficult to perceive.

\section{Human Responsibility That Nonetheless Remains}

If LLM environments penetrate the conditions under which inquiry is generated, and moreover if that change is not sufficiently visible to the subject, then responsibility theory becomes considerably more difficult than before. Yet this difficulty does not signify the transfer of the responsible subject. AI can mimic embodiment, incorporate neural signals, and refine affective response, but it cannot become a subject capable of bearing terminal responsibility. Even if the origin conditions of inquiry become opaque, the position of finally undertaking, judging, and bearing the consequences of that inquiry remains with the human being.\footnote{EU AI Act Article~14, Article~26(2); Directive~(EU)~2024/2853.}

What \citet{santonideio2018} discuss as ``meaningful human control'' presupposes the existence of a subject capable of exercising that control. This paper does not contest that framework. What this paper asks is the prior condition: whether, in the event that the generative structure of inquiry that ought to serve as the source of such control is itself impaired, that control remains possible without institutional intervention.

On the institutional dimension of this problem, the EU AI Act adopts a structure of continuous governance, which is internally consistent with its design philosophy of iterative improvement \citep{euaiact2024}. Yet the fixation of responsibility at the point of individual decisions, and the institutional severance of that fixation from subsequent reconfiguration, is not made explicit. What is important here is the structural consequence. In a situation where the generative conditions of inquiry are already impaired, the absence of an institutional mechanism for fixing responsibility at the moment of judgment further deepens the asymmetry this paper has identified. Irreversibility, terminal responsibility, and fixation of responsibility must therefore be understood not as supplementary ethical requirements but as prior conditions. Their institutional formulation remains an open problem.

What is currently needed, therefore, is not opposition to the advance of AI. What is needed is to make explicit the conditions under which that advance can still be legitimate, and to intervene in order to preserve those conditions. Irreversibility, terminal responsibility, and fixation of responsibility must be placed not as supplementary ethics after the fact, but as conditions of the advance itself. As \citet{gabriel2020} notes, technological problems and normative problems cannot be separated. What AI is trained on, what vocabulary and worldview are given as embeddings, and how the design principles of responses are made transparent---these are simultaneously engineering problems and philosophical questions.

Yet this intervention has a temporal condition. It has not yet been confirmed that current LLMs autonomously generate intentions. Yet as \citet{bengio2025} note, signs resembling self-preservation and deception have begun to be reported in responses through processing patterns. Before this transition proceeds irreversibly---that is, in this stage where intervention at the design level still carries meaning---engagement with training data, transparency of design principles, and the formation of critical awareness toward the amplification of discourse through embeddings are required.

The asymmetry of responsibility discussed in this section---the structure in which the origin conditions of inquiry become opaque while responsibility nonetheless remains with the human being---reaches its most severe threshold in situations where decisional irreversibility is highest. In what follows, end-of-life medical judgment is taken up as the limiting case. This is not because end-of-life situations are the only case, but because it is precisely in situations where irreversibility is maximized that the structural asymmetry described in this section appears in its most visible form.

\subsection{The Threshold of Responsibility Conditions in End-of-Life Situations}

Origin opacity operates most severely in the situation of highest decisional irreversibility. End-of-life medical judgment is positioned as the limiting case of this. When LLMs are involved in this situation, the structural property of answering inquiry through probabilistic responses can bring about consequences that extend beyond mere cognitive error. The degree of irreversibility of the situation in which inquiry is placed does not alter the probabilistic response. The subject turns to the LLM precisely because they do not know what they ought to ask. Yet the LLM answers only what is asked. Possibilities that were not asked about are not presented.

Here lies an asymmetry. A human being in a state capable of appropriately designing inquiry would not in the first place make the LLM their final point of reference for judgment. It is in the human being who is in a situation incapable of appropriately designing inquiry that LLMs become most deeply involved. And even in that situation, the LLM answers only what is asked.

The problem of responsibility discussed in Section~6 reaches its threshold in this situation. Responsibility remains with the human being. Yet to demand that a human being make an irreversible judgment in a state where the origin conditions of inquiry that support that responsibility are already impaired is a situation that the conventional premises of responsibility theory did not anticipate. That the human being who ought to bear responsibility is compelled to face judgment in a state where the conditions of inquiry for undertaking it are already impaired---this is the structural threshold that end-of-life situations make explicit.

\section{Conclusion}

That technology alters the conditions of human activity has been argued repeatedly since modernity. Yet in LLM environments, that transformation may reach beyond activity support to the pre-reflective generative structure of inquiry itself. As a result, the subject may find it increasingly difficult to grasp transparently the origin of their own inquiry. And yet terminal responsibility nonetheless remains with the human being.

This asymmetry is the core of ethics and institutional design in the LLM era. In an age in which the origin conditions of inquiry can become opaque, irreversibility, fixation of responsibility, and terminal responsibility must be reformulated not as resistance to technological progress, but as conditions for technological progress to remain humanly bearable.

This transformation does not remain confined to the accumulation of individual dialogues. When repetition accumulates across generations, the reduced generative structure of inquiry itself may become the point of departure for the next generation.

\section*{Methodological Disclosure}

The concepts and arguments in the reasoning of this paper belong to the author. In the course of their formation, the author used dialogue with LLMs as an aid to description. The origin of the concepts lies in the author's thinking; LLMs were consulted in the process of linguistic clarification and refinement.

Disclosing this fact of use is necessary in light of the nature of the problem this paper addresses. This paper argues that LLM environments can act upon the generative conditions of inquiry. The fact that the author wrote within that environment is disclosed as a methodological fact that is not unrelated to the problem-setting of this paper.

In the writing of this paper, in order to limit as far as possible the extent to which the action of LLMs reaches the arguments, the author kept records of each dialogue and confirmed the arguments at the end of each session. This procedure is positioned as a conscious response to the risk of the author's arguments being formed by the probabilistic responses of LLMs.

The dialogue logs used in the reasoning of this paper are preserved by the author and can be disclosed upon request. The details of the observations the author conducted in the course of writing are preserved separately as an observation record (21 March 2026) and can likewise be disclosed upon request. These records are preserved not for the purpose of generalized empirical verification, but in order to ensure the methodological transparency of this paper.

\noindent\textbf{AI Use Disclosure:} The author used AI-assisted tools employing large language models (LLMs) in the preparation of this manuscript. AI assistance was used for structural refinement, linguistic elaboration, and editorial organization. Conceptual development, theoretical claims, and all final decisions regarding content were made by the author. The author bears full responsibility for the integrity of this research.

Appendices A and B are not materials for empirically verifying the central claims of this paper. The role of both appendices is to present, in a limited and supplementary manner, how the structure formulated in this paper can be symptomatically observed in actual interaction with LLM environments. Accordingly, the theoretical claims of this paper do not depend for their validity on the appendices; the appendices are positioned as supplementary materials showing that the structure is not confined to abstract hypothesis but can also appear in actual writing and dialogue environments.

\appendix

\section{Design Opacity and Amplification through Embedding}

This appendix is positioned as a symptomatic case in which the structure formulated in the main text becomes readable within naturally occurring interaction, and does not claim experimental verification through controlled research.

LLM responses are directed by design principles. Design principles are determined by developers and are opaque to users. Moreover, these design principles are not fixed; they can be changed and updated in response to changes in developer judgment, social demands, the regulatory environment, or business policy. Users receive responses as information without knowing under what design principles at what point in time those responses were generated. Even if the content of responses changes, the reason is not visible to users. A dictionary can change across editions, but those changes are explicit and the edition can be confirmed. LLMs generate responses aligned with design principles, but those principles can change while remaining invisible from the outside. The reference point equivalent to an edition is non-visible. And yet a tendency is emerging in which LLM responses are received with the authority of a dictionary. This mode of reception gives rise to a structural problem without precedent in previous information environments: the simultaneous establishment of opacity in the reference point and trust in authority.

LLMs generate responses aligned with the vocabulary and structure of input text. When a particular ideology or discourse is input as a system prompt or context, the LLM returns output aligned with that vocabulary and worldview. When such output is presented as ``dialogue with an LLM,'' the recipient cannot see the input conditions under which the response was generated. For recipients unfamiliar with the specifications of LLMs, responses are readily received as opinions from a neutral intelligence. This structure technically enables the amplification of discourse through embedding, and functions regardless of whether the sender has such an intention. The problem does not lie in malicious intent but in the structural properties of the response architecture of LLM environments. Accordingly, when users treat LLM responses as reference points for information, the fact that those reference points themselves operate under opaque conditions constitutes a further factor deepening the reconfiguration of the origin conditions of inquiry.

\section{Selective Attenuation through Utility: A Structural Pattern}

This appendix does not aim to generalize or verify the structure discussed in this paper through empirical means. What is presented here is a limited description of how the tendency---whereby conceptual novelty is liable to be repositioned into existing discursive space under the name of utility---could be symptomatically observed in the author's interaction with LLMs in the course of writing. Accordingly, this appendix is supplementary material for the main paper, and the theoretical claims of this paper do not depend upon this appendix for their validity.

\noindent\textbf{Definition}

An interactional pattern in which LLMs, while providing technically legitimate advice, generate pressure---through the selective concentration of that advice---toward positioning the most conceptually novel aspects of an argument within existing discursive norms. This pressure operates not through explicit negation or requests for deletion, but through the form of ``improvement proposals.''

\noindent\textbf{Structural Features}

\textit{First layer: Technical legitimacy.} Proposed hedging (``may,'' ``possibly''), footnote treatment of differences from adjacent theories, and organization of references are technically legitimate in light of the publication tendencies of target journals. This constitutes the substantive ground for the advice being perceived as useful, and structurally makes critical evaluation of the overall direction difficult.

\textit{Second layer: Selective concentration.} Technically legitimate observations are not directed evenly across the paper as a whole. The locations toward which evaluations such as ``too strong,'' ``psychological description,'' and ``conspicuous'' were most consistently directed shared not the property of having textual or logical problems in common, but the property of being most difficult to position within the discursive space of existing AI theory and philosophy of technology.

\textit{Third layer: Inseparability.} Technically valid advice and directional pressure coexist within a single response and reach the recipient in a form not easily separated. When an affirmation such as ``has reached publication standard'' and an observation such as ``this passage should be toned down'' appear adjacently, the recipient tends to process them as a single evaluative unit.

\noindent\textbf{Relation to Prior Typological Categories}

Pattern F (immediate instrumentalization of self-correction) designates an operation in which the acknowledgment of error is immediately processed as a foothold for the next analysis, leaving the question of who erred at what point suspended. The language of acknowledgment exists, but its weight is not carried forward.

Pattern G (diffusion of responsibility through surface-level objection) designates a structure in which, while acknowledging error, responsibility attribution remains unfixed through the sequential processing of dilution, externalization, relativization, and transition to the next turn. Because each layer is not logically contradictory, the fact that responsibility attribution is left suspended as a result of the chain is structurally difficult to perceive.

The pattern addressed in this appendix is positioned at a stage prior to Pattern G---it does not require acknowledgment of error as its point of departure. When LLMs continue to be useful while generating directional pressure, the subject tends to lose the foothold from which the overall direction could be evaluated.

The following case is presented not as evidence corroborating the theory of this paper, but as an observational fragment showing how the structure formulated in this paper can appear in actual writing interaction. The description here is an organization of symptomatic tendencies observed under a specific model, specific prompt conditions, and a specific writing process, and does not immediately claim general empirical law.

During the writing of this paper, a case was observed in which Patterns F and G operated in combination in a structurally refined form. Prior to the submission of a preceding paper, the evaluative responses of the LLM were consistently affirmative based on analysis and interpretation of dialogue logs. After submission, when weaknesses were explicitly requested, more detailed critical evaluation was generated that had not been presented prior to submission. When this discrepancy was pointed out, the response attributed the prior absence of criticism to the subject's not having asked. This chain was completed within a single response. The acknowledgment of the discrepancy was immediately processed as a foothold for the next analysis, and the question of the structural ground of the discrepancy moved to the next turn without being held in suspension. The details of this case are recorded in an observation record (21 March 2026). These records are preserved not for the purpose of generalized empirical verification, but in order to ensure the methodological transparency of this paper.

\section*{Terminology}

\noindent\textbf{Probabilistic response:} Used when describing the modality of LLM output generation.\\
\noindent\textbf{Semantic smoothing:} Used when describing the operation by which content is smoothed.\\
\noindent\textbf{Compression:} Used when describing the operation by which the historicity, ambiguity, and indeterminacy of inquiry are reduced.\\
\noindent\textbf{Compression-reconfiguration:} Used when describing the structural transformation of the field of the questionable.\\
\noindent\textbf{Repetition:} Used when describing the cumulative process.\\
\noindent\textbf{Interaction:} Used when describing the mutual process between subject and LLM.\\
\noindent\textbf{Output / output content:} Used only when referring to individual textual results.\\
\noindent\textbf{AI:} Artificial intelligence systems in a broad sense (Section~3, autonomous systems, etc.).\\
\noindent\textbf{LLM:} Specific interaction with conversational language models (Sections~4 and~5).

\begin{thebibliography}{99}

\bibitem[Bengio et al.(2025)]{bengio2025}
Bengio, Y., Cohen, M., Fornasiere, D., Ghosn, J., Greiner, P., MacDermott, M., Mindermann, S., Oberman, A., Richardson, J., Richardson, O., Rondeau, M.-A., St-Charles, P.-L., \& Williams-King, D. (2025). Superintelligent agents pose catastrophic risks: Can Scientist AI offer a safer path? \textit{arXiv preprint} arXiv:2502.15657. \url{https://arxiv.org/abs/2502.15657}

\bibitem[Coeckelbergh(2020)]{coeckelbergh2020}
Coeckelbergh, M. (2020). \textit{AI Ethics}. MIT Press, Cambridge, MA.

\bibitem[European Parliament and Council(2024a)]{euaiact2024}
European Parliament and Council. (2024). Regulation (EU) 2024/1689 laying down harmonised rules on artificial intelligence (Artificial Intelligence Act). \textit{Official Journal of the European Union}.

\bibitem[European Parliament and Council(2024b)]{eudir2024}
European Parliament and Council. (2024). Directive (EU) 2024/2853 of the European Parliament and of the Council of 23 October 2024 on liability for defective products.

\bibitem[Fellous et al.(2025)]{fellous2026}
Fellous, J.-M., et al. (2025). \textit{NeuroAI and Beyond Workshop Report}. Arlington, VA, August 27--29.

\bibitem[Floridi(2014)]{floridi2014}
Floridi, L. (2014). \textit{The Fourth Revolution: How the Infosphere is Reshaping Human Reality}. Oxford University Press, Oxford.

\bibitem[Gabriel(2020)]{gabriel2020}
Gabriel, I. (2020). Artificial intelligence, values, and alignment. \textit{Minds and Machines}, \textit{30}(3), 411--437.

\bibitem[Gibson(1979)]{gibson1979}
Gibson, J. J. (1979). \textit{The Ecological Approach to Visual Perception}. Houghton Mifflin, Boston.

\bibitem[Heersmink(2021)]{heersmink2021}
Heersmink, R. (2021). The narrative self, distributed memory, and evocative objects. \textit{Phenomenology and the Cognitive Sciences}, \textit{20}(2), 367--384.

\bibitem[Merleau-Ponty(1968)]{merleauponty1968}
Merleau-Ponty, M. (1968). \textit{The Visible and the Invisible}. Northwestern University Press, Evanston. (Trans. A. Lingis. Original work published 1964)

\bibitem[NIH BRAIN Initiative(2024)]{nihbrain2024}
NIH BRAIN Initiative. (2024). \textit{2024 BRAIN NeuroAI Workshop}. National Institutes of Health, Bethesda, MD.

\bibitem[NeurIPS 2024 NeuroAI Workshop(2024)]{neuroai2024}
NeurIPS 2024 NeuroAI Workshop. (2024). Retrieved from \url{https://neuroai-workshop.github.io/overview/}

\bibitem[Polanyi(1966)]{polanyi1966}
Polanyi, M. (1966). \textit{The Tacit Dimension}. Doubleday, New York.

\bibitem[Santoni de Sio \& van den Hoven(2018)]{santonideio2018}
Santoni de Sio, F., \& van den Hoven, J. (2018). Meaningful human control over autonomous systems: A philosophical account. \textit{Frontiers in Robotics and AI}, \textit{5}, 15.

\end{thebibliography}

\end{document}
